\section*{Chapitre \ref{ch:s1:short-term-reversal} --- Retournement de court terme}

\begin{solution}{exo:s1:short-term-reversal:1}
Le poids échangé vaut le double du taux de rotation, donc le coût de point mort vaut $0.303/(2 \times 0.76 \times 252) = 7.9$ points de base par unité échangée. À dix points de base
le coût vaut environ $2 \times 0.75 \times 0.001 \times 252 = 37.8~\%$ par an, soit davantage que le rendement brut.
\end{solution}

\begin{solution}{exo:s1:short-term-reversal:2}
Le rendement enregistré est la fourchette, $0.001 = 10$ points de base, sans changement de valeur. Un signal de retournement fondé sur les prix de transaction repère un gagnant et le vend~; la
transaction du lendemain au cours acheteur «~confirmerait~» le retournement. Il s'agit du rebond bid-ask~: un retournement présent dans les données que personne ne peut négocier, ce qui explique que les signaux soient construits sur
les prix milieux ou les cours de la phase de clôture, et évalués sur les prix auxquels le portefeuille pourrait effectivement être négocié.
\end{solution}

\begin{solution}{exo:s1:short-term-reversal:3}
$0.7 \times 0.02 \times \sqrt{10/25} = 0.00885$~: 88.5 points de base, en plus de la demi-fourchette.
\end{solution}

\begin{solution}{exo:s1:short-term-reversal:4}
$0.04 \times 0.02 \times 2 = 0.0016$, soit 16.0 points de base pour un jour. Une transaction d'un centième du volume quotidien, soit 0.25~million de dollars, coûte
$2 + 10^4 \times 0.7 \times 0.02 \times \sqrt{0.25/25} = 16.0$ points de base~: c'est là que la prévision cesse d'être rentable. La transaction de 10~millions de dollars de l'exercice~3 coûterait
plus de cinq fois l'edge.
\end{solution}

\begin{solution}{exo:s1:short-term-reversal:5}
Le ratio de Sharpe dépend du risque autant que du rendement. Les perdants et les gagnants du portefeuille brut partagent des secteurs et des styles, si bien que le portefeuille porte des paris sur les facteurs qui ajoutent
de la volatilité sans ajouter de rendement~; le portefeuille résiduel en est presque exempt. Le rendement augmente un peu (de 27.0~\% à 30.3~\%), la volatilité baisse beaucoup.
\end{solution}

\begin{solution}{exo:s1:short-term-reversal:6}
Le jour d'une annonce, le mouvement de l'action tient surtout à l'information, et sur un marché présentant une dérive post-résultats, l'information continue de pousser le cours dans le même sens pendant des semaines~: le
portefeuille de retournement achète une action qui continue de baisser. Les dates d'annonce sont programmées et publiées à l'avance, et le filtre n'utilise que le fait de savoir si le jour était un
jour d'annonce, ce qui est connu à la clôture, et non l'effet futur de la surprise.
\end{solution}

\begin{solution}{exo:s1:short-term-reversal:7}
Le ratio de Sharpe net à 300~millions de dollars vaut 0.73. Au-delà de 1~milliard de dollars, le ratio net est statistiquement indiscernable de zéro (0.09, $-0.08$, 0.19 à 1, 3 et 10~milliards de dollars)~:
le portefeuille négocie si peu que son rendement net se réduit à quelques points de base par an, et huit années de données ne suffisent pas à classer des valeurs aussi faibles. Ce qui est robuste, c'est la
baisse, de 1.72 à 100~millions de dollars à environ zéro à 1~milliard de dollars.
\end{solution}

\begin{solution}{exo:s1:short-term-reversal:8}
Le ratio de Sharpe de 13 est calculé avant coûts~; après dix points de base, celui du même portefeuille vaut $-3.43$. La perte à grande taille n'est pas une division par deux~: à 1~milliard de dollars, même le
portefeuille tenant compte des coûts couvre ses coûts et rien de plus. Les bonnes questions portent sur le coût de point mort, le taux de rotation du portefeuille et la courbe de capacité aux
coûts calibrés par la société.
\end{solution}

\begin{solution}{pb:s1:short-term-reversal:1}
\begin{enumerate}
\item Les perdants récents surperforment les gagnants récents sur la courte fenêtre suivante~; un trader qui a besoin d'immédiateté fait bouger le cours, et celui qui absorbe le flux est rémunéré quand
le cours revient.
\item Jegadeesh~: autocorrélation sérielle négative et significative au premier ordre des rendements mensuels, soit 2.49~\% par mois entre les déciles extrêmes sur 1934--1987. Lehmann~: gagnants et perdants hebdomadaires
se retournent la semaine suivante, avec des profits qui résistent aux fourchettes et à des coûts plausibles.
\item Les rendements du retournement se comportent comme ceux de la fourniture de liquidité~: prévisibles à partir du VIX, avec des rendements espérés et des ratios de Sharpe qui bondissent en période de turbulence.
\item L'information modifie la valeur~; seule la part d'un mouvement qui relève de la pression est censée revenir.
\item 0.037, 0.038 et 0.041~; la valeur vraie est 0.044 (la même parmi les 500 titres les plus liquides, 0.039 pour le \omterm{def:s1:short-term-reversal:residual}{retournement ajusté du secteur}).
\item 6.69, 7.24 et 12.96 avant coûts~; $-2.62$, $-2.46$ et $-3.43$ après dix points de base.
\item Il est exempt des expositions au secteur et au style qui ajoutent du risque sans rendement.
\item Le \omterm{def:s1:short-term-reversal:residual}{retournement résiduel} offre des rendements ajustés du risque deux fois plus élevés que le retournement classique et résiste aux coûts dans les grandes capitalisations après 1990~; retirer la part
expliquée par les nouvelles sur les flux de trésorerie a multiplié par quatre le rendement ajusté du risque de la version standard.
\item La règle de Grinold, $\alpha = \mathrm{IC}\,\sigma z$, avec $z$ le résidu rapporté à sa propre volatilité et standardisé, l'IC étant mesuré avant la
évaluation.
\item Rendement espéré moins une pénalité de risque diagonal, moins la demi-fourchette et l'impact en racine carrée, neutre à l'égard des quinze facteurs~; avec un risque diagonal, la transaction de chaque titre
est un pas proximal en forme close, et les multiplicateurs de neutralité proviennent d'une petite itération de la méthode de Newton.
\item 186~\% du capital par jour pour une exposition brute de 2.48~; 246~\%, 738~\% et 2\,293~\% du capital par an en coûts, aux trois tailles.
\item Il écarte les titres le jour de leur annonce~; à 1~milliard de dollars, il transforme $-0.70~\%$ par an en $+0.05~\%$.
\item \textbf{Résultat nommé.} Ratios de Sharpe nets de 1.72 à 100~millions de dollars et de 0.09 à 1~milliard de dollars, avec les coûts intégrés à l'optimiseur, contre $-27.8$ et $-31.7$ pour
le portefeuille qui les ignore.
\item En deçà de 1~milliard de dollars pour ce signal~: l'edge par transaction vaut environ seize points de base pour un titre situé à deux écarts-types, et l'impact en racine carrée atteint
ce niveau à 1~\% du volume quotidien d'un titre.
\item Il a rapporté 10.9~\% par an (ratio de Sharpe 0.93) jusqu'en 1989, et 1.6~\% (0.13) à partir de 1990, avant coûts~; $-5.2~\%$ par an dans les années 2020 à ce jour.
\item C'est un facteur mensuel, pondéré par la capitalisation, sur de grands portefeuilles~; l'effet résiduel quotidien relève d'une autre stratégie, et les rendements du retournement se concentrent dans les crises.
\item Les profits du retournement sont réalisés overnight, avec une continuation au sein de chaque partie de la journée~; le retournement intrajournalier porte sur des déséquilibres de moins d'une heure et du
rebond bid-ask.
\item Avec 50~millions de dollars, le portefeuille résiduel quotidien~; avec 5~milliards de dollars, la surcouche appliquée à un portefeuille plus lent, qui ne coûte presque rien en propre.
\item Le modèle de coûts~: le signe de la stratégie après coûts en dépend.
\item Être rémunéré par une petite commission, plusieurs fois par jour, pour négocier avec des gens pressés, et en conserver plus que ce que coûte la négociation.
\end{enumerate}
\end{solution}

\begin{solution}{iq:s1:short-term-reversal:1}
Les traders qui doivent négocier tout de suite poussent les cours au-delà de la valeur~; le cours revient à mesure que le déséquilibre est absorbé. Le trader de retournement est rémunéré par ces traders pressés,
en tant que teneur de marché à horizon quotidien, et il est le mieux rémunéré lorsque la liquidité se raréfie.
\end{solution}

\begin{solution}{iq:s1:short-term-reversal:2}
Le retournement brut utilise le rendement entier~; le \omterm{def:s1:short-term-reversal:residual}{retournement résiduel} retire ce qu'explique le modèle factoriel. Les mouvements communs relèvent surtout de l'information et ne se retournent pas, et ils donnent au
portefeuille brut des expositions aux facteurs qui ajoutent du risque~; le portefeuille résiduel est plus propre et, sur les données américaines comme ici, affiche un rendement ajusté du risque bien plus élevé.
\end{solution}

\begin{solution}{iq:s1:short-term-reversal:3}
Pas sous la forme d'un portefeuille quotidien naïf~: un coût de point mort de 15 points de base est proche des coûts réels pour la plupart des titres. Il faut demander le modèle de coûts et son calage sur de véritables exécutions, le taux de rotation,
la courbe de capacité avec les coûts intégrés à l'optimiseur, la performance par tranche de liquidité et autour des annonces, ainsi que son comportement dans les crises.
\end{solution}

\begin{solution}{iq:s1:short-term-reversal:4}
Utiliser un modèle factoriel ou à risque diagonal, de sorte que le problème se décompose, ou presque, titre par titre~; les coûts linéaires et en puissance admettent un pas proximal en forme close~; traiter
la neutralité au moyen de quelques multiplicateurs de Lagrange obtenus par la méthode de Newton, ou utiliser la méthode du gradient proximal (FISTA) sur le modèle factoriel. Démarrer à chaud
à partir de la solution de la veille.
\end{solution}

\begin{solution}{iq:s1:short-term-reversal:5}
En période de crise, lorsque les flux qu'il absorbe sont informés ou persistants (ventes forcées qui durent des jours), lorsque les jours d'annonce contaminent le signal, et lorsqu'une
masse de portefeuilles semblables se débouclent. Surveiller les pertes les jours d'annonce, l'exposition aux facteurs, le taux de rotation et les coûts par rapport au modèle, ainsi que la corrélation du portefeuille avec d'autres
portefeuilles de retournement.
\end{solution}

\begin{solution}{iq:s1:short-term-reversal:6}
Le coût par unité échangée d'une transaction de $Q$ dollars vaut $h + \eta\sigma\sqrt{Q/V}$. En l'égalant à $\alpha$, on obtient $Q^* = V\left((\alpha - h)/(\eta\sigma)\right)^2$
pour $\alpha > h$, et aucune transaction sinon. Avec $\alpha$ = 16 points de base, $h$ = 2, $\eta = 0.7$ et $\sigma$ = 2~\%, $Q^* = V/100$.
\end{solution}
